\documentclass[11pt]{article}
\usepackage[T1]{fontenc}
\usepackage{lmodern}
\usepackage[margin=1in]{geometry}
\usepackage{amsmath,amssymb,amsthm}
\usepackage{microtype}
\usepackage[hidelinks]{hyperref}
\hypersetup{pdftitle={Simple negative roots for the rows of the Delannoy matrix square},
  pdfauthor={Xiaojian Zeng}}
\newtheorem{theorem}{Theorem}
\newtheorem{lemma}{Lemma}
\DeclareMathOperator{\sgn}{sgn}
\DeclareMathOperator{\Arg}{Arg}
\title{Simple negative roots for the rows of the Delannoy matrix square}
\author{Xiaojian Zeng\\[3pt]
\small PBC School of Finance, Tsinghua University, 100190\\
\small \texttt{xiaojianzeng@foxmail.com}}
\date{}

\begin{document}
\maketitle
\begin{abstract}
Let $D(n,k)$ count lattice paths from $(0,0)$ to $(n-k,k)$ with steps
$(1,0)$, $(0,1)$ and $(1,1)$. We prove that, for every $n\geq1$, the
$n$th row generating polynomial of the ordinary matrix square $D^2$ is
monic of degree $n$ and has $n$ simple real roots in
$(-4-3\sqrt2,-3+2\sqrt2)$. This establishes the Delannoy component of
Mao and Wang's conjecture on the squares of the Eulerian and Delannoy
triangles, with simplicity and a uniform enclosure. The Eulerian component,
and hence the full conjecture, remains unresolved here. The proof expresses
each row as a divided difference and constructs alternating signs on two
adjacent intervals, using a complex phase on one and a trigonometric
amplitude comparison on the other.
\end{abstract}

\section{Statement and context}

For $n\geq0$ and $0\leq k\leq n$, let $D(n,k)$ count the paths specified
in the abstract, with $D(0,0)=1$ for the empty path. Set $D(n,k)=0$ outside
this range and write
\begin{equation}\label{eq:definitions}
 d_n(y)=\sum_{k=0}^n D(n,k)y^k,
 \qquad
 G_n(x)=\sum_{k=0}^n\sum_{i=k}^n D(n,i)D(i,k)x^k.
\end{equation}
Thus $G_n$ is the row generating polynomial of the ordinary lower
triangular matrix product $D^2$. All sums defining its entries are finite.

Mao and Wang~\cite[Conjecture~4.1, p.~11]{MaoWang2026} conjecture that
the row generating polynomials of both $A^2$ and $D^2$ have only real
nonpositive roots. Here $A(0,0)=1$ and, for $n\geq1$, $A(n,k)$ counts
permutations of $[n]$ with $k-1$ descents. The following theorem concerns
exactly the Delannoy component.

\begin{theorem}\label{thm:main}
With the conventions in~\eqref{eq:definitions}, $G_0=1$. For every integer
$n\geq1$, $G_n$ is monic of degree $n$ and has $n$ distinct real roots,
all in the open interval
\[
 \bigl(-4-3\sqrt2,\,-3+2\sqrt2\bigr).
\]
\end{theorem}

For the ordinary row polynomials $d_n$, Wang, Zheng and
Chen~\cite{WangZhengChen2019} established the root interval
$(-3-2\sqrt2,-3+2\sqrt2)$ and density in its closure. Weighted ordinary
rows, with a corresponding simplicity and interval theorem, are treated by
Chen, Zhang and Pan~\cite[Theorem~1]{ChenZhangPan2026}. These results supply
the ordinary-row background; Theorem~\ref{thm:main} treats the composed rows
$G_n=\sum_i D(n,i)d_i$.

Mao and Wang also derive the recurrence
\begin{equation}\label{eq:fourterm}
 G_n=(x+3)G_{n-1}+(x-1)(G_{n-2}+G_{n-3})+xG_{n-4}
 \qquad(n\geq4),
\end{equation}
with $G_0=1$, $G_1=x+2$, $G_2=x^2+6x+5$ and
$G_3=x^3+10x^2+25x+12$~\cite[Section~4]{MaoWang2026}.
Alternating-sign arguments are a standard tool for real-rootedness;
Liu and Wang~\cite[Theorem~2.1 and its proof]{LiuWang2007} develop them
in an interlacing framework. Our argument constructs the sample points
directly for each $G_n$. It does not require interlacing of consecutive
rows or a general preservation property of the Delannoy transformation.

The contribution is the all-degree Delannoy-square assertion together with
the stated simplicity and enclosure. No equivalent assertion was found in
the primary material inspected through September 14, 2026. This is a bounded
literature assessment: access to the 2019 paper was limited to its publisher
abstract and bibliographic record, and we make no exhaustive priority claim.

\section{A divided difference}

Last-step decomposition gives
\begin{equation}\label{eq:ordinary}
 d_0(y)=1,\quad d_1(y)=1+y,\quad
 d_n(y)=(1+y)d_{n-1}(y)+y d_{n-2}(y)\quad(n\geq2).
\end{equation}
The characteristic-root method for this recurrence also appears in the
weighted setting of~\cite[Section~2]{ChenZhangPan2026}. We give the
identities needed here explicitly. If $w,z$ are the distinct roots of
$t^2-(1+y)t-y=0$, the initial values and recurrence imply
\begin{equation}\label{eq:quadratic}
 d_n(y)=\frac{w^{n+1}-z^{n+1}}{w-z}.
\end{equation}
Let $u,v$ be the distinct roots of $t^2-(1+x)t-x=0$. Reordering the
finite sums in~\eqref{eq:definitions} and using~\eqref{eq:quadratic} yields
\begin{equation}\label{eq:divided}
 G_n(x)=\sum_{i=0}^n D(n,i)d_i(x)
       =\frac{u d_n(u)-v d_n(v)}{u-v}.
\end{equation}
The factors $u$ and $v$ arise from the exponent $i+1$ in the formula
for $d_i(x)$.

For the rest of the proof fix $n\geq1$ and put
\begin{equation}\label{eq:constants}
 \begin{gathered}
 N=n+1,\qquad a=-4-3\sqrt2,\quad b=-3-2\sqrt2,\quad c=-3+2\sqrt2,\\
 u_0=-1-\sqrt2,\qquad
 \theta_0=\arccos\frac{1+u_0}{2\sqrt{-u_0}}.
 \end{gathered}
\end{equation}
Then $a<b<c<0$ and $\pi/2<\theta_0<\pi$.

\section{Two families of alternating signs}

\subsection{The interval (b,c)}

For $b<x<c$ choose the outer roots as
\[
 u=\frac{1+x+i\sqrt{-(x^2+6x+1)}}2,\qquad v=\overline u.
\]
They satisfy $\operatorname{Im}u>0$ and $|u+1|^2=2$. As $x$ decreases
from $c$ to $b$, $u$ traverses the upper semicircle from $\sqrt2-1$ to
$u_0$. In particular $\operatorname{Re}u>-3$, so
\[
 \operatorname{Im}(u^2+6u+1)
 =2\operatorname{Im}u\,(\operatorname{Re}u+3)>0.
\]
Let $s=\sqrt{u^2+6u+1}$ have positive real and imaginary parts and set
\begin{equation}\label{eq:inner}
 w=\frac{1+u+s}2,\qquad z=\frac{1+u-s}2,\qquad B=\frac{u}{s}.
\end{equation}
These functions are continuous on the arc and $s=w-z\ne0$.
Moreover $\operatorname{Im}w>0$. The quadratic relations give
$z=(1-w)/(1+w)$ and hence
\[
 \operatorname{Im}z=-\frac{2\operatorname{Im}w}{|1+w|^2}<0.
\]
Here $1+w\ne0$, and $wz=-u\ne0$.

\begin{lemma}\label{lem:dominance}
For every $x\in(b,c)$ the roots in~\eqref{eq:inner} satisfy $|w|>|z|$.
\end{lemma}
\begin{proof}
If $|w|=|z|$, then $w/z$ has modulus one and
\[
 \frac{(w+z)^2}{wz}=2+\frac wz+\frac zw\in[0,4].
\]
The same expression is $-u-2-1/u$, whose imaginary part equals
$\operatorname{Im}u\,(|u|^{-2}-1)$. Its vanishing forces $|u|=1$.
Together with $|u+1|^2=2$ and $\operatorname{Im}u>0$, this implies $u=i$.
But $-i-2-1/i=-2$, a contradiction. Thus equality never occurs on the
connected arc. At the endpoint $x=c$ the limiting $u$ is positive, and
the limiting real roots satisfy $w>0>z$ and $w+z=1+u>0$. Consequently
$|w|>|z|$ near that endpoint, and continuity proves the assertion throughout.
\end{proof}

Use the continuous arguments $\Arg u\in(0,\pi)$,
$\Arg s\in(0,\pi/2)$ and $\Arg w\in(0,\pi)$, and define
\[
 \Phi(x)=N\Arg w+\Arg u-\Arg s.
\]
As $x\to c^-$, the three quantities $u,s,w$ tend to positive real
numbers, so $\Phi(x)\to0$. As $x\to b^+$,
\[
 u\to u_0,\qquad s\to i\sqrt{2+4\sqrt2},\qquad
 w\to\frac{-\sqrt2+i\sqrt{2+4\sqrt2}}2
       =\sqrt{-u_0}\,e^{i\theta_0}.
\]
Therefore
\begin{equation}\label{eq:limits}
 \Phi(c^-)=0,\qquad \Phi(b^+)=N\theta_0+\frac\pi2.
\end{equation}
Since $d_n$ has real coefficients, equations~\eqref{eq:quadratic}
and~\eqref{eq:divided} give
\begin{equation}\label{eq:phase}
 G_n(x)=\frac{\operatorname{Im}\{B(w^N-z^N)\}}{\operatorname{Im}u}.
\end{equation}
At a point where $\Phi(x)=(j+\tfrac12)\pi$, the imaginary part of
$Bw^N$ is $(-1)^j|Bw^N|$. Lemma~\ref{lem:dominance} makes
$|Bz^N|<|Bw^N|$, and thus
\begin{equation}\label{eq:sign}
 \sgn G_n(x)=(-1)^j.
\end{equation}

For every integer $j$ with $0\leq j<N\theta_0/\pi$, the level
$(j+\tfrac12)\pi$ lies strictly between the limits in~\eqref{eq:limits}.
The points realizing these levels can be chosen in strictly decreasing
order of $x$. Indeed, parametrize the closed arc from $c$ to $b$ and
extend $\Phi$ by its endpoint limits. For each level take its first hitting
time. Continuity ensures existence in the open arc, and reaching a higher
level from zero requires passing every lower level strictly earlier.
No monotonicity of $\Phi$ is needed.

\subsection{The interval (a,b)}

For $u\in(b,u_0)$ put
\[
 h(u)=\frac{u(u-1)}{u+1},\qquad x=h(u),\qquad v=\frac{1-u}{1+u}.
\]
The identities
\[
 h(b)=a,\qquad h(u_0)=b,\qquad h'(u)=1-\frac{2}{(u+1)^2}>0
\]
show that $h$ maps $(b,u_0)$ increasingly onto $(a,b)$. Also
$u+v=1+x$, $uv=-x$ and
\begin{equation}\label{eq:ordering}
 b<u<u_0<v<-\sqrt2<-1.
\end{equation}
Thus $u,v$ are distinct outer roots in~\eqref{eq:divided}, and both lie
in $(b,c)$.

On $(b,c)$ define $\theta(t)\in(0,\pi)$ by
\[
 \cos\theta(t)=\frac{1+t}{2\sqrt{-t}}.
\]
The derivative of the right side is $(1-t)/(4(-t)^{3/2})>0$;
hence $\theta$ decreases strictly with $t$. In the range~\eqref{eq:ordering},
\[
 \frac\pi2<\theta(v)<\theta(u)<\pi.
\]
As $u$ decreases from $u_0$ to $b$, $\theta(u)$ increases from
$\theta_0$ to $\pi$. Applying~\eqref{eq:quadratic} to the conjugate
roots $\sqrt{-t}\,e^{\pm i\theta(t)}$ gives
\[
 d_n(t)=(-t)^{n/2}\frac{\sin(N\theta(t))}{\sin\theta(t)}.
\]
Put $M(t)=(-t)^{n/2+1}/\sin\theta(t)$. Then $M(u)>M(v)>0$,
because $-u>-v>0$ and
$0<\sin\theta(u)<\sin\theta(v)$. Equation~\eqref{eq:divided} becomes
\begin{equation}\label{eq:amplitude}
 G_n(h(u))=
 \frac{-M(u)\sin(N\theta(u))+M(v)\sin(N\theta(v))}{u-v}.
\end{equation}
For each integer $j$ satisfying
\begin{equation}\label{eq:leftindices}
 \frac{N\theta_0}{\pi}<j+\frac12<N,
\end{equation}
there is exactly one $u\in(b,u_0)$ with
$\theta(u)=(j+\tfrac12)\pi/N$. At this point the first term in the
numerator of~\eqref{eq:amplitude} strictly dominates the second.
Since $u-v<0$, its sign is again $\sgn G_n(h(u))=(-1)^j$.
The corresponding $x=h(u)$ decrease as $j$ increases.

\subsection{Joining the samples and counting roots}

Let $q=N\theta_0/\pi$ and $m=\lceil q\rceil$. As $0<q<N$,
we have $1\leq m\leq N$. For $j=0,\ldots,m-1$, use the samples
in $(b,c)$: the required inequality $j<q$ holds even if $q$ is an
integer. For $j=m,\ldots,N-1$, use the samples in $(a,b)$, since
$q<j+\tfrac12<N$. If $m=N$, this second group is empty.

There are exactly $N=n+1$ samples, in the order
\[
 c>x_0>x_1>\cdots>x_n>a,\qquad \sgn G_n(x_j)=(-1)^j.
\]
The intermediate value theorem gives a root in each of the $n$ disjoint
intervals $(x_{j+1},x_j)$. The unique vertical path gives $D(n,n)=1$,
so $(D^2)(n,n)=1$ and $G_n$ is monic of degree $n$. Thus these are
all its roots and each is simple. They lie in $(a,c)$, as asserted.
The join $x=b$ needs no sign evaluation and may lie inside one of these
root intervals. Finally $G_0=1$ follows from the empty path. This proves
Theorem~\ref{thm:main}.

\section{Exact checks and remaining questions}

The argument above covers every degree without computation. The
supplementary checks used Python~3.10.12, SymPy~1.14.0~\cite{SymPy2017}, and
mpmath~1.3.0 for sample discovery. Exact rational root isolation confirmed
the assertion for every $1\leq n\leq24$. For
\[
 n\in\{1,2,3,4,5,8,16,31,64,128,257\},
\]
the stored certificates contain all $n+1$ alternating sample signs,
530 signs in total. Floating-point phase inversion located candidate
samples; subsequent sign, order and enclosure decisions used exact
integer or rational arithmetic. Separate implementations rebuilt the
path counts and the ordinary matrix products and replayed these signs.
These checks cover only the stated finite degrees.

For reproducibility, write $G_n(x)=\sum_{k=0}^n g_kx^k$. At a rational
sample $p/r$, with $r>0$, its sign is the sign of the integer
\begin{equation}\label{eq:certificate}
 H_n(p,r)=\sum_{k=0}^n g_kp^kr^{n-k}.
\end{equation}
Given the coefficients, the $n+1$ sample checks form a multipoint
polynomial evaluation task. Zeng~\cite{Zeng2026Evaluation}
studies its generic characteristic-zero complexity in relation to
derivative Bézout inversion and interpolation. Here direct integer
evaluation suffices; we make no fast-complexity claim. As discussed
in~\cite[Section~2.1]{Zeng2026Evaluation}, a bit-cost analysis would
also have to account for the sizes of the integers, beyond counting
arithmetic operations.

The enclosure can be checked at each rational sample without approximating
$\sqrt2$. For a real rational $t$, the strict inequalities $a<t<c$ are
equivalent to
\[
 \bigl(t+4\geq0\ \text{or}\ (t+4)^2<18\bigr),\qquad
 \bigl(t+3\leq0\ \text{or}\ (t+3)^2<8\bigr).
\]
Together with strict sample order and the signs in~\eqref{eq:certificate},
these tests reproduce the finite root certificates. Each certificate is
specified by its coefficient vector, a common positive denominator,
and the ordered integer numerators of the sample points.

Theorem~\ref{thm:main} does not establish consecutive strict interlacing,
optimality of either endpoint, or a general Delannoy preservation theorem.
The Eulerian-square component of the original conjecture remains
unresolved here. Each of these questions requires an additional argument.

\paragraph{Computational and AI assistance.}
OpenAI Codex assisted with proof development, symbolic and exact-arithmetic
verification programs, source comparison, a separate-context mathematical
review, and manuscript preparation. That review was an additional model
pass, not external peer review. No proof assistant was used. The mathematical
justification is the explicit argument above; the finite computations and
AI review are supporting checks.

\bibliographystyle{plain}
\bibliography{references}
\end{document}
